\section*{अध्याय \ref{ch:b1:structure-spectroscopy} --- संरचना के लिए स्पेक्ट्रमिकी}

\begin{solution}{exo:b1:structure-spectroscopy:1}
$\varepsilon = A/(\ell c) = 0.64/(1.00 \times \num{2.0e-5}) =
\qty{3.2e4}{L/(mol.cm)}$। तीन गुनी सांद्रता: $A = 1.92$, यदि
उस \omterm{def:b1:structure-spectroscopy:absorbance}{अवशोषकता} पर भी नियम लागू रहता हो।
\end{solution}

\begin{solution}{exo:b1:structure-spectroscopy:2}
\ce{C6H6}: $D = (12 + 2 - 6)/2 = 4$, बेंज़ीन (एक वलय और तीन $\pi$ आबंध)।
\ce{C4H8O2}: 1, एथिल एथेनोएट या ब्यूटेनोइक अम्ल। \ce{C3H7NO}: $(6 + 2 + 1
- 7)/2 = 1$, प्रोपेनऐमाइड। \ce{C2H3Cl}: 1, क्लोरोएथीन। \ce{C10H14O}:
$(20 + 2 - 14)/2 = 4$: कार्वोन में एक वलय, दो C=C और एक C=O है।
\end{solution}

\begin{solution}{exo:b1:structure-spectroscopy:3}
प्रोपेनोन: एक संकेत (6H)। एथेनोइक अम्ल: दो (3H, 1H)। 2-मेथिलप्रोपेन:
दो (9H, 1H)। 1,2-डाइक्लोरोएथेन: एक (4H)।
\end{solution}

\begin{solution}{exo:b1:structure-spectroscopy:4}
$\delta = 1460/400 = \qty{3.65}{ppm}$; \qty{90}{MHz} पर संकेत संदर्भ से
$3.65 \times 90 = \qty{329}{Hz}$ दूर होता है। \omterm{def:b1:structure-spectroscopy:coupling}{युग्मन स्थिरांक}
हर्ट्ज़ में वही रहते हैं जबकि विस्थापनों के अंतर क्षेत्र के साथ बढ़ते हैं:
जो \omterm{def:b1:structure-spectroscopy:coupling}{बहुक} निम्न क्षेत्र पर अतिव्यापित होते हैं, वे उच्च क्षेत्र पर पृथक हो जाते हैं।
\end{solution}

\begin{solution}{exo:b1:structure-spectroscopy:5}
\ce{CH3}: त्रिक (3H); केंद्रीय \ce{CH2}: पाँच पड़ोसी, षष्टक
1\,:\,5\,:\,10\,:\,10\,:\,5\,:\,1 (2H); \ce{CH2Br}: त्रिक (2H), सबसे अधिक
विपरिरक्षित।
\end{solution}

\begin{solution}{exo:b1:structure-spectroscopy:6}
एथेनॉल: O--H (3666) और कोई C=O नहीं; प्रोपेनोन: C=O (1738) और कोई O--H नहीं;
एथेनोइक अम्ल: दोनों, O--H (3582) और C=O (1788)।
\end{solution}

\begin{solution}{exo:b1:structure-spectroscopy:7}
एथिल एथेनोएट: 4.1 के पास चतुष्क (2H), 2.0 के पास एकक (3H), 1.3 ppm के पास त्रिक
(3H)। मेथिल प्रोपेनोएट: \ce{OCH3} का एकक (3H), ऑक्सीजन द्वारा विपरिरक्षित
(3.7 ppm के पास), \ce{CH2C=O} का चतुष्क (2H) (2.3 के पास), 1.1 के पास त्रिक
(3H)। पहले में चतुष्क सबसे अधिक विपरिरक्षित संकेत है और
दूसरे में नहीं: 4 ppm के पास चतुष्क का अर्थ है \ce{O-CH2CH3}।
\end{solution}

\begin{solution}{exo:b1:structure-spectroscopy:8}
रेखाएँ \qty{0.080}{ppm} दूर हैं, अर्थात् $0.080 \times 90 = \qty{7.2}{Hz}$:
\omterm{def:b1:structure-spectroscopy:coupling}{युग्मन स्थिरांक}। त्रिक का अंतराल भी यही है: \ce{CH2} और
\ce{CH3} एक-दूसरे से युग्मित हैं।
\end{solution}

\begin{solution}{exo:b1:structure-spectroscopy:9}
एक C=O पट्ट और एक ही संकेत: प्रोपेनोन, \ce{CH3COCH3}। ऐल्डिहाइड
समावयवी, प्रोपेनल \ce{CH3CH2CHO}, \ce{CHO} समूह का C--H तनन कंपन
और उसका अत्यधिक विपरिरक्षित प्रोटॉन दर्शाता है।
\end{solution}

\begin{solution}{exo:b1:structure-spectroscopy:10}
पहला समुच्चय रेखा को $n_1 + 1$ रेखाओं में विभक्त करता है; इनमें से हर एक को
दूसरा समुच्चय $n_2 + 1$ में विभक्त करता है: $(n_1 + 1)(n_2 + 1)$ रेखाएँ। यदि $J_1 =
J_2$, तो किसी रेखा की स्थिति केवल $k_1 + k_2$ पर निर्भर करती है, जो $n_1 + n_2 +
1$ मान लेता है: $n + 1$ नियम, $n = n_1 + n_2$ के साथ। 1-ब्रोमोप्रोपेन का केंद्रीय \ce{CH2}:
सामान्यतः $4 \times 3 = 12$ रेखाएँ, और दोनों
स्थिरांक बराबर होने पर छह (एक षष्टक)।
\end{solution}

\begin{solution}{exo:b1:structure-spectroscopy:11}
$D = 0$; एक ऐल्कोहॉल (O--H पट्ट, कोई C=O नहीं)। द्विक (6H): एक CH पर दो तुल्य
\ce{CH3}; \omterm{def:b1:structure-spectroscopy:coupling}{बहुक} (1H): वही CH; द्विक (2H): CH के पास और OH वहन करने वाला
\ce{CH2}; चौड़ा एकक: OH।
2-मेथिलप्रोपेन-1-ऑल, \ce{(CH3)2CH-CH2-OH}।
\end{solution}

\begin{solution}{exo:b1:structure-spectroscopy:12}
हर स्पीशीज़ स्वतंत्र रूप से अवशोषित करती है: $A = A_1 + A_2 = \ell(\varepsilon_1c_1
+ \varepsilon_2c_2)$। दो तरंगदैर्घ्यों पर, चारों गुणांक ज्ञात होने पर,
दो रैखिक समीकरण $c_1$ और $c_2$ देते हैं।
\end{solution}

\begin{solution}{pb:b1:structure-spectroscopy:1}
\textbf{1.} $0.5690 \times 12.0/44.0 = \qty{0.1552}{g}$।
\textbf{2.} $0.2328 \times 2.0/18.0 = \qty{0.02587}{g}$।
\textbf{3.} $0.2500 - 0.1552 - 0.0259 = \qty{0.0689}{g}$।
\textbf{4.} C \qty{62.1}{\percent}, H \qty{10.3}{\percent}, O
\qty{27.6}{\percent}।
\textbf{5.} C $0.1552/12.0 = 0.01293$, H $0.02587$, O $0.0689/16.0 =
0.00431$ mol: अनुपात $3.00 : 6.00 : 1$, \ce{C3H6O}।
\textbf{6.} $M = \rho RT/p = 3740 \times 8.314 \times 373.15/(1.000 \times 10^5) =
\qty{116}{g/mol}$ (घनत्व \unit{g/m^3} में)।
\textbf{7.} $116/58 = 2$: \ce{C6H12O2}, $D = (12 + 2 - 12)/2 = 1$।
\textbf{8.} एक C=O तनन कंपन।
\textbf{9.} कोई भी O--H: न अम्ल, न ऐल्कोहॉल।
\textbf{10.} एक C--O तनन कंपन, प्रबल: एस्टर का C--O।
\textbf{11.} \ce{CH2} और \ce{CH3} समूहों के C--H तनन कंपन।
\textbf{12.} एक $\pi$ आबंध, जिसे C=O ने ले लिया है; प्रबल C--O के साथ और
O--H के बिना, एक एस्टर। अम्ल को O--H चाहिए; डाइऑल में C=O नहीं होता। ईथर समूह
वहन करने वाले कीटोन का C=O प्रोपेनोन के \qty{1738}{cm^{-1}} के पास होता,
जबकि प्रेक्षित \qty{1754}{cm^{-1}} एथिल एथेनोएट के
\qty{1764}{cm^{-1}} के निकट है; भाग III का NMR एस्टर की
पुष्टि करता है।
\textbf{13.} पाँच; $2 + 2 + 2 + 3 + 3 = 12$, बारह हाइड्रोजन।
\textbf{14.} एक \ce{CH2}, तीन पड़ोसियों (एक \ce{CH3}) के साथ, ऑक्सीजन से
आबंधित (विपरिरक्षित): \ce{-O-CH2-CH3}।
\textbf{15.} उसी एथिल समूह का \ce{CH3}, दो पड़ोसियों के साथ; दोनों
संकेत एक-दूसरे से युग्मित हैं, इसलिए एक ही $J$।
\textbf{16.} दो पड़ोसियों वाला \ce{CH2}, C=O के पास।
\textbf{17.} पाँच पड़ोसियों वाला \ce{CH2}: एक \ce{CH2} और एक
\ce{CH3} के बीच।
\textbf{18.} दो पड़ोसियों वाला \ce{CH3}, क्रियात्मक समूह से दूर।
\textbf{19.} \ce{CH3CH2O-} और \ce{-CH2CH2CH3}।
\textbf{20.} 4.13 ppm वाले प्रोटॉन ऑक्सीजन से आबंधित कार्बन पर हैं,
2.28 ppm वाले C=O के पास: दोनों इलेक्ट्रॉन-आकर्षी।
\textbf{21.} \ce{CH3CH2-O-CO-CH2CH2CH3} या \ce{CH3CH2-CO-O-CH2CH2CH3}। केवल
पहला तीन पड़ोसियों वाले \ce{CH2} को ऑक्सीजन पर रखता है (4.13 ppm पर
चतुष्क)।
\textbf{22.} प्रोपिल प्रोपेनोएट दूसरा है: उसका \ce{OCH2} एक
त्रिक होता और चतुष्क 2.3 ppm के पास होता।
\textbf{23.} ब्यूटिल एथेनोएट, \ce{CH3CO-O-C4H9}, 2.0 ppm के पास एक एकक (3H)
दर्शाता और कोई चतुष्क नहीं।
\textbf{24.} एथिल ब्यूटेनोएट, \ce{CH3CH2CH2-CO-O-CH2CH3}।
\textbf{25.} अवरक्त: C=O 1754, C--O 1182, C--H 2982, कोई O--H नहीं। NMR: \ce{OCH2}
4.13 (चतुष्क), \ce{CH2CO} 2.28 (त्रिक), केंद्रीय \ce{CH2} 1.66 (षष्टक), एस्टर का
\ce{CH3} 1.25 (त्रिक), शृंखला का \ce{CH3} 0.95 (त्रिक), \omterm{def:b1:structure-spectroscopy:integration}{समाकलन} 2, 2, 2, 3, 3।
\textbf{26.} \textbf{\qty{116}{g/mol}}: एथिल ब्यूटेनोएट, \ce{C6H12O2}।
\end{solution}
